\section{Cyber Access \& Governance}
\sectionkicker{AUTHORIZED ACCESS}
\begin{wideflow}
{\Large\bfseries Cyberモデルの提供は、許可の設計から始まる}\par
\smallskip
{\color{SurveyMuted}GPT-5.6-Cyber / Daybreak Redの動きは、一般向けモデル公開ではなく、authorized access、partner distribution、safeguardの境界を読む材料として整理する。}\par
\end{wideflow}

OpenAIの確認資料は、GPT-5.6-CyberをDaybreak Red経由で、authorized vulnerability researchとsecurity testingに利用できるものとして記述している\cite{daybreak-defense}。ここで示されるのは「誰でも使える新モデル」ではなく、目的を限定したprogram accessである。記事では、この主体と用途の帰属を保ったまま扱う。

同じDaybreakでも、アクセス経路は一つではない。別のOpenAI資料はDaybreak modelsをAmazon Bedrock経由でenterprise security workflowsに利用できるものとして記し、さらにapproved Daybreak partnersがfrontier cyber modelsをauthorizedでgovernedなcybersecurity servicesに使う形も記述する\cite{daybreak-aws,daybreak-partners}。Bedrockやpartnerの記録はDaybreakの重複launchではなく、distributionと統治の補助線として置く。

\begin{themeoverview}[三つの範囲を分ける]
\begin{tabularx}{\linewidth}{@{}>{\bfseries}p{0.25\linewidth}X@{}}
Model scope & GPT-5.6-Cyber / Daybreak Redとして記述されたモデルの範囲。\\
Access scope & vulnerability research、security testing、enterprise workflow、approved partner serviceという目的・経路の範囲。\\
Safeguard boundary & 確認資料はauthorized / governedという方向を示すが、具体的なsafeguardsと一般のmodel/API availabilityまでは閉じていない。
\end{tabularx}
\end{themeoverview}

この三分割は、利用可否を一語で答えないためのものだ。program accessは、目的、主体、経路、統治条件を含む。したがって、その説明を一般提供の宣言へ拡張しない。Bedrockでのavailabilityも、記録が示すaccess routeを超えて推測しない。

\begin{claimboundary}[次に見るべき境界]
確認資料が示すauthorized / governedという説明を超えて、具体的な safeguards、全モデルの範囲、一般のAPI availabilityを断定しない。次の観測点は、program accessと一般提供を区別できるモデル範囲、Bedrockでのavailability境界、そして safeguards の具体化である。
\end{claimboundary}

この章での問いは、「そのモデルは使えるか」から「誰が、どの目的で、どの経路を通り、どの統治条件の下で使えるか」へ移る。アクセスの制度設計そのものが、モデルの実用面を規定する。
